\documentclass{homework}

\title{04}
\date{\DTMdate{2025-12-06}}

\usepackage{fancyvrb}
\usepackage{hyperref}
\fvset{listparameters=\setlength{\topsep}{0pt}\setlength{\partopsep}{0pt}}

\begin{document}
\maketitle

\textbf{Reminders}

\begin{itemize}
  \item I will give you feedback on GitLab and Discord.
  Check them after receiving your (preliminary) grade.
  \item Read through task descriptions carefully.
  \item You may request deadline extensions until 03.12.25 23:59 EET.
  \item You can submit up to one week after the deadline with a 50\% point penalty.
  0 afterwards.
  \item Grace clause:
  \begin{itemize}
    \item If you submit your assignment by the deadline, you have the opportunity to fix (some) mistakes I point out to gain back up to half of the lost points.
    \item You cannot get any points back for a task that is entirely missing or broken, e.g. code that does not run at all or does nothing.
    If your code does not pass the assertions/test script (if available), include that info in your report.
    Otherwise I will consider your code to be broken, i.e. no grace.
    Test your code!
    \item The grace clause does not apply if you requested an extension or missed the original deadline.
  \end{itemize}
  \item You may share hints with each-other, but no answers/code.
  Ask questions in the server if they could be useful for everyone.
  \item You may---and I encourage you to---ask me questions (preferably during practices).
  \item You may use AI as help, but not to solve the tasks themselves.
  \begin{itemize}
    \item Acceptable: e.g. checking your understanding, asking how to loop in Python
    \item Not acceptable: e.g. generating any code logic or cryptographic operations
    \item If you use AI in any form, include in your report what you used it for
    \item I recommend that you ask in the server/me instead of using AI
  \end{itemize}
  \item If I catch you plagiarising, copying, or using AI for theory answers or code logic, I will report you to the faculty and fail you for the course.
  \begin{itemize}
    \item If during the midterm/exam you completely fumble stuff that you got points for in homework tasks, you will have to orally explain why it is so.
    \item If you do not convince me of homework authorship, I will report you for cheating and fail you.
  \end{itemize}
\end{itemize}

\newpage

\begin{center}
  Theory tasks
\end{center}

\textbf{Additional instructions}

\begin{itemize}
  \item The report file should preferably use the Markdown format, and make use of Markdown syntax for headings, code blocks, etc.
  I shall accept no formats other than markdown and PDF.
  \item Call the file \texttt{report.md} (or \texttt{report.pdf} if you must use PDF) and include it in the \texttt{hw4} folder root.
  \item Your homework must contain the following information at the start:
\begin{Verbatim}
- Name: <your full name>
- Discord username: <your discord username>
- Feedback language: <EE or EN>
- Homework took me:
  - Theory: <time taken for the theory>
  - Practical: <time taken for the practical tasks>
\end{Verbatim}
\end{itemize}

\begin{task}{Proof systems \& zero-knowledge}
  For this task, answer the following questions using complete sentences:
  \begin{itemize}
    \item Why should a proof system be sound?

    \item Are interactive proofs typically transferable?
    That is, should an observer who observes an interactive proof but trusts neither the prover nor the verifier be convinced by the proof?

    \item What mechanism can be used to transform an interactive proof into a non-interactive one?
    Do the security guarantees remain the same?

    \item Consider an interactive proof system that provides $256$ bits of security.
    Could you maintain this security level by using SHA-256 in the non-interactive version of the proof?
    Why yes/why not?
  \end{itemize}

  Then, describe a use case where a zero-knowledge proof would be useful in practice, which is not one of the examples discussed in class (age range, citizenship of member states, online voting).

  If you used online materials to help you answer these questions, references in your answers.
\end{task}

\newpage

\begin{task}{TLS}
  For this task, answer the following questions:
  \begin{itemize}
    \item Describe three aspects in which TLS 1.3 improves on TLS 1.2.
    Explain in your own words why you think these are important improvements.

    \item Pick a website that uses HTTPS, and answer the following (include screenshots):
    \begin{itemize}
      \item What is the website link?
      \item Is the certificate domain validated, organisation validated, or extended validated?
      How do you know?
      \item What is the validity period of the certificate?
      \item Where is the OCSP responder or CRL?
      \item What is the physical address of the root CA?
    \end{itemize}

    \item Why do we need certificate transparency?
  \end{itemize}

  If you used online materials to help you answer these questions, please include references in your answers.
\end{task}

\begin{task}{eMRTD}
  For this task, answer the following questions based on \textit{\href{https://blog.trailofbits.com/2025/10/31/the-cryptography-behind-electronic-passports/}{this blog post on electronic passports}}.
  You need not (but I recommend that you) read through the entire article.
  \begin{itemize}
    \item What types of biometric data are included in your passport?
    \item Is this data protected?
    If so, how and by whom can it be extracted?
    \item How is the authenticity of passports verified?
    \item Why should you not leave your passport unattended?
  \end{itemize}

  There may be a question of similar style about passports/the EE ID card in the exam.
\end{task}

\begin{task}{A taste of the exam}
  For this task, answer at least 10 of following questions using complete sentences.
  I encourage you to answer more, but 10 is the minimum.

  Be succinct and precise: no word soup!
  In the exam, if you are not able to give a clear answer, I will consider that you do not know the answer.

  \begin{enumerate}
    \item Does cryptography protect the availability of information?
    Why yes/why not?
    \item Why do we need encryption?
    \item Why is IND-CPA important?
    \item Why do MACs not offer non-repudiation?
    \item Why is public key crypto not safe against unbounded adversaries?
    \item Does second pre-image resistance imply collision resistance?
    \item Why is using a cryptographic hash function for hashing passwords a bad idea?
    \item What is important when protecting passwords?
    \item What is one role of a key derivation function?
    \item What are block cipher modes?
    \item What is the advantage of AES-GCM over AES-CTR (or of ChaCha-Poly over ChaCha)?
    \item What do digital signatures provide?
    \item What is typically needed on top of digital signatures for trust?
    \item Why do we need eIDAS?
    \item How is the authenticity of a server checked when you visit a website?
    \item How do devices/browsers trust root CAs?
    \item What is the role of a CA?
    \item What is needed for the OTP to be unconditionally secure?
    \item Why is the OTP unconditionally secure?
    \item What are some advantages of OCSP over CRLs?
    \item What are the two core properties of a commitment scheme?
    \item Why does an identification protocol require \enquote{freshness}?
  \end{enumerate}

  These questions you should be able to answer on your own.
  If you need to look up the answers, then that is a good indicator that you will need to prepare for the exam and go through materials and/or recordings again.
\end{task}

\begin{task}{Your thoughts (optional)}
  This is an \emph{optional} task, but I'd appreciate it if you at least answer the first question.
  You can also choose to answer only some of the questions.
  It's not anonymous (clearly) but this is because I might want to ask you questions if you say something particularly intriguing.
  It's fine to say negative things about the course and/or me, I only ask for the answer to be constructive.

  \begin{itemize}
    \item What's your mood? (select up to 2)
    \begin{enumerate}
      \item You killed crypto for me, I hope to never deal with it!
      \item I liked crypto more before this class.
      \item It was ok, but that's the last I'll deal with it.
      \item It was interesting, but that's the last I'll deal with it.
      \item I like crypto more after this class.
      \item It was interesting, I may pursue crypto in the future.
      \item I aspire to be the next Daniel Bernstein.
    \end{enumerate}
    \item Do you feel like this course has been useful for you? (it's okay to say no)
    \item Do you feel like this course has been fair? (again it's ok to say no)
    \item What has/have been the hardest topic(s)?
    \item What would you have liked to be different?
    \item What would you have liked to see more/less of?
    \item Any additional thoughts/suggestions for future iterations of this class?
  \end{itemize}
\end{task}

\newpage
\setcounter{task}{0}

\begin{center}
  Practical tasks
\end{center}

\begin{task}{Signing your commits}
  You must set up a commit signing key on TalTech's GitLab.
  I recommend that you use an SSH key for signing commits, but you can also use GPG or S/MIME.
  I will not include a reference here: you should be able to find one yourself (GitLab and GitHub both have articles for this).

  For the task, you must push a signed commit into your homework repository.
  This commit should include the public key of the key you use for signing: call it simply \texttt{key}.
  The key should not be a binary blob: it should open in a text editor (e.g. PEM).
  Your task is considered complete if the GUI shows that your commit is verified.

  \begin{figure}[h!]
    \center
    \includegraphics[width=\textwidth]{../images/gitlab_verified.png}
    \caption{Verified commit on GitLab.}
  \end{figure}

  Include the full commit link in your report.
\end{task}

\begin{task}{SSH \& SCP}
  For this task, you must SSH into \texttt{enos.itcollege.ee} using an SSH key.

  If you have not yet set up an SSH key, you will need to get it onto the server, e.g. with \texttt{ssh-copy-id}\footnote{\url{https://www.ssh.com/academy/ssh/copy-id\#copy-the-key-to-a-server}}.
  The connection will ask you for your account password on your connection: this is your UniID password.

  Once you can SSH into Enos using an SSH key, take a screenshot of the log-in flow.
  If your SSH connection does not prompt you for a password, explain in writing how you have protected your key (e.g. password in SSH agent, SSH key on Yubikey).
  For example, my SSH key is a non-resident key on my Yubikey, and the SSH key password is automatically loaded by my SSH agent.

  \begin{figure}[h!]
    \center
    \includegraphics[width=\textwidth]{../images/enos_ssh.png}
  \end{figure}

  Once in the server, create an empty directory called \texttt{ics0036} in your \texttt{public\_html} folder.
  Then, close the connection with \texttt{exit}.

  Upload the screenshot that you took of the log-in flow into the \texttt{\textasciitilde/public\_html/ics0036} folder on the server using SCP (secure copy)\footnotemark{}.
  Take a screenshot of this process as well, and add it also to the folder on the server.
  \footnotetext{\url{https://linuxize.com/post/how-to-use-scp-command-to-securely-transfer-files/}}

  Finally, include \emph{both} screenshots also in your report.
  To verify whether your screenshots are correctly on the server, visit \url{https://enos.itcollege.ee/~takraa/ics0036} but with your own student username instead of \texttt{takraa}.
\end{task}
\end{document}
